\documentclass[12pt]{article}
\usepackage[utf8]{inputenc}
\usepackage[T1]{fontenc}
\usepackage{hyperref}
\title{Tabla Language Resources for K-12 Curriculum Development}
\author{Compiled by Claude Sonnet 5 (claude-sonnet-5) for Andr\'as Kornai}
\date{}
\begin{document}
\maketitle

\section{Introduction}
Tabla (ISO 639-3 tnm; Glottocode tabl1243) is a Sentani-branch language of the East Bird's Head--Sentani family. It is spoken on the coast of Tanahmerah Bay near Jayapura, Papua, Indonesia. About 3,800 speakers are cited from 1990. SABDA, drawing on Ethnologue, rates it as moribund. Dialects named are Yokari, Tepera and Yewena-Yongsu. Joshua Project has a ``Tepera, Tanahmerah'' profile that places the people near Tanah Merah town on the Digul River and calls the language an isolate. That conflicts with the other sources and should not be relied on.

\section{Primary Educational Resources}
None found. No primers, textbooks or school programmes for Tabla appeared in the search results.

\section{Linguistic and Reference Materials}
\begin{itemize}
\item Tabla language (encyclopedia overview): \url{https://en.wikipedia.org/wiki/Tabla_language}
\item Tabla people: \url{https://en.wikipedia.org/wiki/Tabla_people}
\item Verbix endangered-languages page on a different language, Tanahmerah (Sumeri, tcm). It is useful only to avoid confusion: \url{https://docs.verbix.com/EndangeredLanguages/Tanahmerah}
\end{itemize}
No grammar or dictionary was found.

\section{Religious and Community Resources}
\begin{itemize}
\item SABDA language entry for tnm: \url{https://bakat.sabda.org/604-tnm-001.html}
\item Joshua Project, Tepera, Tanahmerah (classification disputed, see above): \url{https://joshuaproject.net/people_groups/15347}
\end{itemize}
No scripture translation was seen in the results.

\section{Digital and Audio Resources}
None found.

\section{Curriculum Development Context}
Almost nothing is documented online. Any curriculum would depend on new documentation work with elders, since the language may be spoken mainly by older adults. Indonesian is the likely language of instruction. Fieldwork or archive searches (not done here) are needed before planning.
\end{document}